\subsection{The auxiliary abelian tower and its strict cohomology classes}
\label{sec:support-auxiliary-tower}

The auxiliary field must meet three requirements at once: its completions
at the distinguished prime must be large, the relevant ideal class ranks
must remain bounded as its cyclotomic layer grows, and two prescribed
strict cohomology classes must survive restriction.  The following
construction supplies these requirements for the same fields and also
constructs the Galois quotient used in
Theorem~\ref{thm:support-finite-end}.  Its bound is uniform in the
cyclotomic layer, with the initial arithmetic data fixed.

Write $G_{\mathbb Q}=\operatorname{Gal}(\overline{\mathbb Q}/\mathbb Q)$.
For a place $v$ of $\overline{\mathbb Q}$ over a rational prime, let
$D_v$ and $I_v$ denote its decomposition and inertia subgroups.  Let
$\ell$ be an odd prime, $k$ a finite field of characteristic $\ell$,
$\varepsilon:G_{\mathbb Q}\to\mathbb Z_\ell^\times$ the cyclotomic
character, and $\omega$ its reduction to $k^\times$.  The original
character $\Psi$ takes values in the units of a local coefficient
ring with residue field $k$; write $\psi=\overline\Psi$.

\begin{theorem}[The actual auxiliary tower]
\label{thm:support-auxiliary-tower}
Let $\eta$ be the character of an imaginary quadratic field
$K=\mathbb Q(\sqrt{-D})$, where $D>0$.  Fix a positive integer
$M$ divisible by $2\ell$, containing every prime factor of $D$.
Assume the following conditions.
\begin{enumerate}
\item The finite characters $\eta$ and $\psi$ factor through
  $\operatorname{Gal}(\mathbb Q(\mu_{M^a})/\mathbb Q)$ for some
  $a\geq1$.  The character $\eta$ is unramified at $\ell$, and
  $\psi$ is even.  For some $j\geq0$, its restriction to inertia
  at $\ell$ is $\omega^j$.
\item The inclusion $\ker\omega\cap\ker\psi\subseteq\ker\eta$
  fails.  At the fixed place $B$ above $\ell$ there is
  $t\in D_B$ with $\eta(t)=1$ and $\psi(t)^2\ne1$.
\item There is a fixed abelian extension $F_b/\mathbb Q$ of
  $\ell$-power degree, unramified at $\ell$.  If $p\ne\ell$,
  $v\mid p$, and $s\in I_v\cap G_{F_b}$ has $\psi(s)=1$,
  then $\Psi(s)=1$.
\end{enumerate}
For every prescribed positive integer $Q$, one can choose primes
$q,w$, a real quadratic field $F_0$, and the degree-$\ell$
subfield $F_w$ of $\mathbb Q(\mu_w)$, independently of $m$, such
that, with $\mathbb Q_m$ the degree-$\ell^m$ subfield of
$\mathbb Q(\mu_{\ell^{m+1}})$,
\[
 F_m^+=F_0F_wF_b\mathbb Q_m,\qquad
 L_m=F_m^+K,\qquad U_m=G_{F_m^+},                         \tag{AT1}
\]
the following assertions hold for every $m\geq0$.
\begin{enumerate}
\item $F_m^+$ is totally real and abelian over $\mathbb Q$;
  $L_m$ is abelian and CM, with maximal real subfield $F_m^+$.
  Its quadratic extension $L_m/F_m^+$ is unramified at every
  finite place and split at every place above $M$.
\item At every $v\mid\ell$,
  $[F_{m,v}^+:\mathbb Q_\ell]\geq\ell^m$,
  $\omega(I_v\cap U_m)\ne\{1\}$, and there is
  $t_v\in D_v\cap U_m$ with $\psi(t_v)^2\ne1$.
\item If $E_m$ is the fixed field of
  $U_m\cap\ker\eta\cap\ker\psi\cap\ker\omega$, then
  there is an integer $C_0$, independent of $m$, such that
  $\dim_{\mathbb F_\ell}\operatorname{Cl}(E_m)/
       \ell\operatorname{Cl}(E_m)\leq C_0$.
\item For each $b\in\{1,2\}$, put $\kappa_b=\psi^{-b}\eta$.
  There is a continuous cocycle $z_b:G_{\mathbb Q}\to k(\kappa_b)$
  whose restriction to $U_m$ is not a coboundary, whose
  restriction to every $D_v$ with $v\mid\ell$ is a coboundary,
  and which satisfies
  \[
   z_b(s)=0\quad\text{if }p\ne\ell,\ v\mid p,
   \ s\in I_v\cap U_m,\ \eta(s)=\psi(s)=1.               \tag{AT2}
  \]
  In fact the same cocycle $z_b$ can be chosen for all $m$.
\end{enumerate}
The prime $w$ is prime to $M$, satisfies $w\equiv1\pmod{\ell Q}$,
and splits completely in the fields cut out by $\eta$, $\psi$
and $\omega$.  In particular $m$ may be chosen, after $C_0$,
so that $\ell^m>3C_0+7$.
\end{theorem}

\begin{proof}
Choose a prime $q>4D+M$ with $q\equiv-1\pmod{4M}$, and set
$F_0=\mathbb Q(\sqrt{Dq})$.  Dirichlet's theorem applies because
$-1$ is coprime to $4M$.  Since $q\nmid D$, this is a real
quadratic field.  Put $K'=\mathbb Q(\sqrt{-q})$.
For an odd $p\mid M$, the congruence $-q\equiv1\pmod p$
and Hensel's lemma show that $K'$ is split at $p$.
For $p=2$, the congruence is modulo $8$, since $2\mid M$;
an odd $2$-adic unit is a square exactly when it is $1$ modulo
$8$.  Thus $K'$ is split at every prime dividing $M$.
The character of $F_0$ is the product of the characters of
$K$ and $K'$.  It therefore restricts to $\eta$ on every
$D_v$ above $M$.  It is ramified at $q$, where $\eta$ is
unramified.  In particular the character of $F_0$ cannot be
a character unramified outside $M$.

Choose $w>q$ prime with
$w\equiv1\pmod{\ell Q M^a}$.  This again uses only Dirichlet's
theorem.  Its Frobenius acts trivially on all the indicated
roots of unity, so the splitting assertions hold.  The cyclic
group $(\mathbb Z/w\mathbb Z)^\times$ has a unique quotient
of order $\ell$, giving $F_w$.  It is totally real, unramified
away from $w$, and totally ramified of degree $\ell$ at $w$.
The fields $F_b$ and $\mathbb Q_m$ are also totally real,
because they are Galois of odd degree over $\mathbb Q$.
This proves the asserted properties of $F_m^+$.

Complex conjugation fixes $F_m^+$ and acts nontrivially on
$K$.  Hence $[L_m:F_m^+]=2$ and $F_m^+$ is its maximal real
subfield.  At a prime dividing $M$, the equality of the two
quadratic local characters says that $K$ is contained in the
local field supplied by $F_0$, unless both are already split.
Thus $L_m/F_m^+$ is split there.  Away from $M$, $K/\mathbb Q$
is unramified; unramifiedness is preserved by local base change.
This proves the first assertion without any restriction on
ramification in the other factors of (AT1).

The completion of $\mathbb Q_m$ at $\ell$ is totally ramified
of degree $\ell^m$, which gives the local degree lower bound.
The fields $F_0,F_w,F_b$ are unramified at $\ell$.  On inertia
at $\ell$, the mod-$\ell^{m+1}$ cyclotomic character is
surjective.  Choose an element with cyclotomic value $-1$.
It fixes the three unramified factors and fixes $\mathbb Q_m$,
because $(-1)^{\ell-1}=1$.  Its value under $\omega$ is $-1$,
which is not $1$.  This proves the second cyclotomic assertion;
in particular none of these completions contains $\mu_\ell$.

Conjugate the given $t\in D_B$ to a place $v\mid\ell$.
Characters are conjugation invariant, so it still satisfies
$\eta(t)=1$ and $\psi(t)^2\ne1$.  It fixes $F_0$ by the
local equality of quadratic characters.  If $[F_b:\mathbb Q]
=\ell^h$, then $t^{\ell^{h+m+1}}$ fixes each of
$F_w,F_b,\mathbb Q_m$ and hence belongs to $D_v\cap U_m$.
The order of $\psi(t)$ is prime to $\ell$, so taking this
power does not make its square equal to $1$.

For the class-rank bound, put
\[
 E=F_0F_wF_bK K_\psi\mathbb Q(\mu_\ell),
 \qquad K_\psi=\overline{\mathbb Q}^{\ker\psi}.
\]
This is one fixed abelian number field.  The definition of
$\mathbb Q_m$ gives the exact equality
\[
 E_m=E(\mu_{\ell^{m+1}}).                                \tag{AT3}
\]
Indeed, on $(\mathbb Z/\ell^{m+1}\mathbb Z)^\times$, the
kernel of the $(\ell-1)$st power intersects the kernel of
reduction modulo $\ell$ trivially.  Therefore adjoining
$\mathbb Q_m$ and $\mu_\ell$ is exactly adjoining
$\mu_{\ell^{m+1}}$.

Ferrero--Washington says that the unramified Iwasawa module
of the cyclotomic $\mathbb Z_\ell$-extension of this abelian
$E$ has $\mu=0$\cite{AuxTowerFW}.  The structure theorem
then makes it finitely generated over $\mathbb Z_\ell$:
the elementary factors are quotients by distinguished
polynomials, with finite kernels and cokernels, and there
are no factors $\mathbb Z_\ell[[T]]/(\ell^r)$.
After a fixed initial level the primes above $\ell$ are
totally ramified in the cyclotomic tower.  Norm maps on the
$\ell$-primary ideal class groups are then surjective.
To see the last assertion, the cokernel of the norm for a
cyclic degree-$\ell$ step is the Galois group of the
intersection of the step with the Hilbert class field of
the lower field.  A nontrivial degree-$\ell$ step ramified
at $\ell$ has trivial such intersection.  The inverse limit
therefore surjects onto every sufficiently late class group.
A fixed number of $\mathbb Z_\ell$-generators bounds all their
$\ell$-ranks.  Enlarge this number to include the finitely
many earlier fields.  This proves the assertion with $C_0$.
The tower structure and the norm argument are also described
in \cite{AuxTowerIwasawa}, Sections 2.4 and 3.2.

It remains to construct the two classes.  Fix $b\in\{1,2\}$
and write $\kappa=\kappa_b$.  It is odd, because
\[
 \psi(\sigma)=1,\qquad \eta(\sigma)=-1,
 \qquad \kappa(\sigma)=-1                               \tag{AT4}
\]
for complex conjugation $\sigma$.  This calculation uses
the parity of the actual character $\psi$, including its
cyclotomic factor.  By hypothesis there is
$s\in\ker\omega\cap\ker\psi$ with $\eta(s)=-1$; hence
$\kappa(s)=-1\ne\omega(s)$, so $\kappa\ne\omega$.
At $\ell$, the element with $\eta(t)=1$ and
$\psi(t)^2\ne1$ makes $\kappa|_{D_B}$ nontrivial for both
$b=1$ and $b=2$.  The same holds at every conjugate place.

Let $S$ consist of $\ell,w$ and the primes at which $\kappa$
is ramified.  On $H^1(\mathbb Q_p,k(\kappa))$ prescribe the
zero condition for $p\in S\setminus\{w\}$, the full
condition at $w$, and the unramified condition outside $S$.
At infinity use the zero condition.  Denote this Selmer
group by $H^1_{\mathcal L}$.  Finite Poitou--Tate duality
and the local and global Euler formulas give
\[
 \dim_k H^1_{\mathcal L}(\mathbb Q,k(\kappa))
 -\dim_k H^1_{\mathcal L^\perp}
          (\mathbb Q,k(\kappa^{-1}\omega))
 =1.                                                       \tag{AT5}
\]
Here every term can be checked separately.  Both global
$H^0$ terms vanish, since $\kappa\ne1,\omega$.  At each
finite prime in $S\setminus\{w\}$ the local $H^0$ term
vanishes: at $\ell$ this was just proved, and at the other
primes the character is ramified.  At $w$ both $\kappa$ and
$\omega$ are trivial; local duality and the Euler formula
at $w\ne\ell$ give $\dim H^1=2$ and $\dim H^0=1$.
Thus its contribution is $2-1=1$.  The real-place term is
zero by (AT4).  These are the full terms of the
Greenberg--Wiles formula derived from the duality statements
in \cite{AuxTowerDuality}; no sign or local term is omitted.
Since the second dimension in (AT5) is nonnegative, choose
a nonzero class and a continuous cocycle representative $z_b$.

The restriction of this class to every $U_m$ is nonzero.
Indeed $U_m$ is normal, and the element $t_v$ constructed
above has $\kappa(t_v)\ne1$.  Hence $k(\kappa)^{U_m}=0$.
Inflation--restriction gives an injection
\[
 H^1(G_{\mathbb Q},k(\kappa))\longrightarrow
                       H^1(U_m,k(\kappa)).               \tag{AT6}
\]
This does not require the index of $U_m$ to be prime to
$\ell$.  The local zero conditions give the required
coboundaries at every place above $\ell$.

For (AT2), first let $p\ne w,\ell$.  If $\kappa$ is
ramified at $p$, the local class is zero, so on inertia
$z_b(s)=(\kappa(s)-1)c$ for some $c$.  If it is unramified,
the unramified local condition gives zero on inertia.
In either case $\eta(s)=\psi(s)=1$ implies $z_b(s)=0$.
At $w$, the local character $\kappa$ is trivial, so $z_b$
is a homomorphism on $D_w$.  Its restriction to inertia
kills wild inertia, which is pro-$w$, and factors through
the tame pro-$\ell$ quotient $\mathbb Z_\ell$.  Its target
has exponent $\ell$, so it kills $\ell\mathbb Z_\ell$.
The extension $F_w/\mathbb Q$ has inertia degree $\ell$ at
$w$, and $I_w\cap G_{F_w}$ maps exactly to this subgroup.
Since $U_m\subseteq G_{F_w}$, (AT2) follows also at $w$.
The choice of $z_b$ preceded $m$, as asserted.
\end{proof}

\begin{corollary}[The prescribed Galois quotient and local maps]
\label{cor:support-auxiliary-galois-quotient}
For every field in Theorem~\ref{thm:support-auxiliary-tower},
let $N_m$ be the closed normal subgroup of $U_m$ generated
by all elements
\[
 s\in I_v\cap U_m,\qquad v\nmid\ell,\qquad
                   \eta(s)=\psi(s)=1,
 \qquad \Gamma_m=U_m/N_m.                                \tag{AT7}
\]
Then the quotient map $p:U_m\to\Gamma_m$ and its actual
decomposition groups give precisely the Galois quotient,
markings, local maps and local degrees required in
Theorem~\ref{thm:support-finite-end}.  The two classes above
descend to nonzero strict classes on $\Gamma_m$.
\end{corollary}

\begin{proof}
The characters $\eta,\psi,\varepsilon$ kill $N_m$.
For $\eta$ and $\psi$ this follows from the generators; for
$\varepsilon$ it follows from the absence of cyclotomic
ramification away from $\ell$.  Their descended characters
have exactly their original values after composition with
$p$.  This is a continuous surjection of profinite groups,
with the literal kernel (AT7).

Take one place above each prime of $F_m^+$ over $\ell$.
Each may be written $\gamma B$ with $\eta(\gamma)=1$:
start with any such representative, and if necessary
multiply $\gamma$ on the left by complex conjugation in
$U_m$.  This changes neither its $U_m$-orbit nor the
prime of $F_m^+$, and reverses the sign of $\eta(\gamma)$.
Let $G_v=U_m\cap D_{\gamma B}$, give it its actual profinite
topology, take the local inclusion $j_v:G_v\hookrightarrow U_m$,
and define its local quotient map as $p\circ j_v$.  Its
inertia group is $U_m\cap I_{\gamma B}$.  Splitting at
$\ell$ puts its image in $\ker\eta$.  Its prescribed
integer degree is exactly
\[
 [D_{\gamma B}:D_{\gamma B}\cap U_m]
                         =[F_{m,v}^+:\mathbb Q_\ell].    \tag{AT8}
\]
The chosen representatives give the required exhaustive,
unique list of $U_m$-orbits of places.  Each local inclusion
is continuous and injective, with the specified image.

Take $c=p(\sigma)$, where $\sigma$ is the original complex
conjugation, and $g_0=p(t_B)$ with the local $t_B$ in the
proof of the theorem.  Then $c^2=1$, $\eta(c)=-1$,
$\psi(c)=1$, $\eta(g_0)=1$ and $\psi(g_0)^2\ne1$.
For every real embedding of $F_m^+$ take its actual
complex conjugation; their images give the remaining
archimedean markings.  The number of these embeddings is
$[F_m^+:\mathbb Q]=\sum_{v\mid\ell}[F_{m,v}^+:\mathbb Q_\ell]$.
The required original condition $\Psi=1$ on the indicated
inertia elements follows from $U_m\subseteq G_{F_b}$ and
the last hypothesis of the theorem.  All other splitting,
unramifiedness and local cyclotomic assertions were proved
for these very fields in the theorem.

Finally $z_b$ vanishes on the generators of $N_m$ by (AT2).
As $\kappa_b=1$ there, the cocycle identity shows that the zero locus
on $\ker\kappa_b$ is a normal subgroup of $U_m$.  It is
closed because $k$ is finite and $z_b$ is continuous.
Thus $z_b|_{U_m}$ descends through $p$.  A descended
coboundary would inflate to a coboundary on $U_m$,
contrary to (AT6), and its local classes are zero by
the original local conditions.  Choosing local primitives
gives nonzero classes in the prescribed relative
cohomology groups as well: a relative coboundary would
in particular make the global class a coboundary.
\end{proof}

The class-rank constant $C_0$ may depend on $\ell$, $q$, $w$,
$F_b$, $\eta$ and $\psi$.  No uniform bound over changing
initial data, and no bound for the orders of growing Selmer
groups, follows from this construction.  The construction
discharges the actual auxiliary field, strict-class and
Galois-quotient requirements.
Theorem~\ref{thm:support-auxiliary-frame} and
Proposition~\ref{prop:support-auxiliary-seed} supply the
finite-coefficient comparisons and seed conditions.
Global integral rigidity and the weighted Selmer estimate
still require their separate arguments.
